\begin{abstract}
在数字化转型与人工智能热潮持续升温的背景下，越来越多 IPO 企业在招股说明书和交易所问询回复中强调自身的智能化、平台化与数字化能力。然而，企业对 AI 与数字化的高频表述究竟反映了真实转型能力，还是更多体现为概念化包装和叙事修辞，仍缺乏细化识别。围绕这一问题，本文拟以 A 股 IPO 企业中明确提及 AI 或数字化内容的企业为研究对象，结合招股说明书与问询回复文本，构建企业 AI 相关披露真实性综合指数，并检验其对上市后经营表现的影响。

具体而言，本文从提及强度、场景具体性和表述审慎性三个维度刻画企业 AI 披露的真实性，并保留叙事包装性这一反向面向，形成“真实性-包装性”双面识别框架。方法上，本文采用“关键词识别-相关段落抽取-LLM 辅助标注-人工复核-指数构建-基础回归验证”的技术路线，其中 LLM 主要承担摘要、标签建议与初步评分任务，关键标签由人工复核确认。经营表现主要从 ROA、营业收入增长率和总资产周转率三个维度衡量，并进一步构建综合经营表现指标进行补充分析。

本文的潜在贡献主要体现在三个方面：第一，将研究窗口前移至 IPO 阶段，讨论企业 AI 叙事究竟是真实能力信号还是包装性表达；第二，提出一套适用于公开披露文本的轻量化真实性识别框架；第三，为投资者识别概念包装、为监管部门理解 IPO 企业相关披露质量、为企业规范数字化叙事提供参考。

\noindent\textbf{关键词：}IPO；人工智能披露；数字化转型；叙事包装；经营表现
\end{abstract}

\begin{center}
\textbf{Abstract}
\end{center}

Against the continued rise of digital transformation and artificial intelligence narratives, an increasing number of Chinese IPO firms emphasize intelligent upgrading, platform construction, and digital capability in prospectuses and exchange inquiry responses. Yet frequent references to AI do not necessarily indicate substantive transformation. They may also reflect concept packaging and impression management. This study focuses on A-share IPO firms that explicitly mention AI- or digital-related content and develops a composite index of AI disclosure authenticity based on prospectuses and inquiry-response texts, then examines whether the index predicts post-IPO operating performance.

The proposed framework evaluates disclosure authenticity from three dimensions: mention intensity, scenario specificity, and statement cautiousness, while retaining narrative packaging as a reverse dimension. Methodologically, the study follows a practical path of keyword identification, passage extraction, LLM-assisted coding, manual verification, index construction, and baseline regression analysis. Post-IPO performance is primarily measured by ROA, revenue growth, and total asset turnover, with a composite operating-performance index used in supplementary tests.

This study may contribute in three ways. First, it shifts the discussion of digital narratives to the IPO stage, where disclosure is dense and regulatory scrutiny is strong. Second, it offers a lightweight but transparent measurement framework that is feasible for undergraduate-level empirical execution. Third, it provides implications for investors, regulators, and firms in distinguishing substantive digital transformation from symbolic narrative packaging.
